% O6. Class H's takeover share in Proposition 2.
% The named statement is the sentence that part (ii) of Proposition 2 gains in Section 4.1, after the display of the
% ownership condition (sections/04_equilibrium_results.tex:39). The paragraph after it is Step 4 of the proof of
% Proposition 2 in Appendix A.3, in place of appendix/A3_proposition2.tex:75-86.
% source: referee report W2, "Shared Candidate 3" (the odds (e_H/e_L)(W_H/W_L) and the exact finite condition);
% W2, "Do the same workarounds cover Proposition 2" (outbidding at F_1 forces it at F_0); PLAN.md Part 5, O6.
% Definitions used (Section 2.4): W_j(F) = E_j[F(theta) 1{theta >= p}], the probability that a prepared challenger
% of class j beats the incumbent F; class H outbids F when W_H(F) = 1; class H's takeover share is the probability
% that the challenger is of class H, given that the challenger acquires the target.

\begin{named}{Proposition~\ref{prop:ccc}(ii), the takeover share}
Under the hypotheses of Proposition~\ref{prop:ccc}, suppose that class \(H\) outbids \(F_1\), \(W_H(F_1)=1\). Then
class \(H\) outbids \(F_0\) as well, high-class ownership rises from \(\pi\rho\) to \(\pi[\rho+(1-\rho)(b+\lambda)]\),
and class \(H\)'s takeover share rises strictly if \(W_L(F_0)>0\) and equals one against both incumbents if
\(W_L(F_0)=0\).
\end{named}

\emph{Step 4: entry, ownership and the takeover share at \(F_1\).} Class \(H\) places the order \(1\) with
probability \(f_H(1)=b+\lambda\), through a liquidity purchase or an informed purchase. Class \(L\) places it with
probability \(f_L(1)=b\). Only the order \(1\) recruits the high-cost challenger. Hence
\[
e_H=\rho+(1-\rho)(b+\lambda),\qquad e_L=\rho+(1-\rho)b,\qquad
\Ent_1=\pi e_H+(1-\pi)e_L=\rho+(1-\rho)(b+\pi\lambda)>\rho ,
\]
since \(b>0\); also \(e_H-e_L=(1-\rho)\lambda>0\). Preparation depends on the order and the cost, and both are
independent of \((\theta,R)\) given the class. A prepared challenger of class \(j\) acquires the target when
\(\theta>R\), with probability \(W_j(F_1)\). Write \(\Own_{j,i}\) for the probability that the challenger is of class
\(j\) and acquires the target at \(F_i\), with \(\pi_H=\pi\) and \(\pi_L=1-\pi\). Then
\(\Own_{j,1}=\pi_je_jW_j(F_1)\), and Step~1 gives \(\Own_{j,0}=\pi_j\rho\,W_j(F_0)\). Division by \(\pi>0\) shows
that high-class ownership rises exactly when \([\rho+(1-\rho)(b+\lambda)]\,W_H(F_1)>\rho\,W_H(F_0)\).

The challenger acquires the target with positive probability against both incumbents. By (H1),
\(B_1(M)>c_H>0\), so \(g_{j,1}>0\) for some class \(j\); a class whose gross profit is positive beats the incumbent
with positive probability, so \(W_j(F_1)>0\). Since \(F_1\) strengthens \(F_0\), \(F_1\le F_0\) pointwise, and
\(W_j(F_0)\ge W_j(F_1)\) for both classes. Class \(H\)'s takeover share at \(F_i\) is
\(\Own_{H,i}/(\Own_{H,i}+\Own_{L,i})\). For \(x,y,x',y'\ge0\) with \(x+y>0\) and \(x'+y'>0\), the inequality
\(x/(x+y)>x'/(x'+y')\) holds exactly when \(xy'>x'y\). With \(x=\Own_{H,1}\), \(y=\Own_{L,1}\), \(x'=\Own_{H,0}\) and
\(y'=\Own_{L,0}\), and after division by \(\pi(1-\pi)\rho>0\), the share rises from \(F_0\) to \(F_1\) exactly when
\begin{equation}\label{eq:a3-share}
e_H\,W_H(F_1)\,W_L(F_0)>e_L\,W_H(F_0)\,W_L(F_1).
\end{equation}

Let class \(H\) outbid \(F_1\). Then \(1\ge W_H(F_0)\ge W_H(F_1)=1\), so class \(H\) outbids \(F_0\), and
\(\Own_{H,1}=\pi e_H>\pi\rho=\Own_{H,0}\). Condition~\eqref{eq:a3-share} reads \(e_H\,W_L(F_0)>e_L\,W_L(F_1)\). If
\(W_L(F_0)>0\), then \(e_H\,W_L(F_0)>e_L\,W_L(F_0)\ge e_L\,W_L(F_1)\), and the share rises strictly. If
\(W_L(F_0)=0\), then \(W_L(F_1)=0\), so \(\Own_{L,0}=\Own_{L,1}=0\), and the share equals one against both
incumbents. With outbidding, the stronger incumbent changes only the low class's chance to win, and entry rises more
in class \(H\) than in class \(L\), so the takeovers that go to the challenger shift toward class \(H\).
